\documentclass[12pt]{article}
\usepackage{../../ficha}

\begin{document}

\ficha{Vetores, Produto Escalar e Geometria no Espaço}

\instrucoes{Considera um referencial ortonormado. Apresenta os cálculos e simplifica os resultados.}

\secao{Vetores no Plano}{Considera $\vec u = (2, -1)$, $\vec v = (-3, 4)$ e os pontos $A(1, 2)$ e $B(5, -1)$. Calcula:}
\begin{exercicios}
  \item $\vec u + \vec v$
  \item $2\vec u - \vec v$
  \item $\|\vec v\|$
  \item As coordenadas de $\vect{AB}$ e $\|\vect{AB}\|$.
  \item Um vetor unitário com a direção e sentido de $\vec v$.
  \item O ponto $C = A + 2\vec u$.
  \item Determina $k$ de modo que $(k, 6)$ seja colinear com $\vec u$.
  \item Escreve $\vec w = (7, -6)$ como combinação linear de $\vec u$ e $\vec v$.
\end{exercicios}

\secao{Vetores no Espaço}{Considera os pontos $A(1, 0, 2)$, $B(3, -1, 4)$ e $C(0, 2, 1)$.}
\begin{exercicios}
  \item Determina as coordenadas de $\vect{AB}$ e de $\vect{AC}$.
  \item Calcula $\|\vect{AB}\|$.
  \item Calcula a distância entre $B$ e $C$.
  \item Determina as coordenadas do ponto médio de $[BC]$.
  \item Determina as coordenadas de $D$ de modo que $ABDC$ seja um paralelogramo.
  \item Escreve a equação da superfície esférica de centro $A$ e raio 3.
  \item Indica o centro e o raio da superfície esférica $x^2 + y^2 + z^2 - 2x + 4z - 4 = 0$.
  \item Escreve a equação do plano mediador de $[AB]$.
\end{exercicios}

\secao{Produto Escalar}{Calcula o produto escalar e, quando pedido, o ângulo entre os vetores.}
\begin{exercicios}
  \item $\vec u = (3, 4)$ e $\vec v = (4, -3)$. São perpendiculares?
  \item $\vec u = (1, 2)$ e $\vec v = (3, 1)$. Determina o ângulo entre eles.
  \item $\vec u = (1, 0, 1)$ e $\vec v = (0, 1, 1)$. Determina o ângulo entre eles.
  \item $\vec u = (2, -1, 3)$ e $\vec v = (1, 5, 1)$.
  \item Sabendo que $\|\vec u\| = 4$, $\|\vec v\| = 3$ e o ângulo entre eles é $60^\circ$, calcula $\vec u \cdot \vec v$.
  \item Determina $k$ de modo que $(k, 2, -1)$ e $(3, k, 4)$ sejam perpendiculares.
  \item Num triângulo de vértices $A(0,0)$, $B(4,0)$ e $C(1,3)$, determina a amplitude do ângulo $B\hat{A}C$.
  \item Determina a equação da reta tangente à circunferência $x^2 + y^2 = 25$ no ponto $(3, 4)$.
\end{exercicios}

\secao{Retas e Planos no Espaço}{Resolve.}
\begin{exercicios}
  \item Escreve uma equação vetorial da reta que passa em $A(1, 2, 3)$ com vetor diretor $(2, -1, 1)$.
  \item Escreve uma equação vetorial da reta que passa em $A(0, 1, -1)$ e $B(2, 3, 1)$.
  \item O ponto $(5, 0, 5)$ pertence à reta do exercício 1?
  \item Escreve uma equação do plano que passa em $A(1, 0, 2)$ e é perpendicular ao vetor $(2, -1, 3)$.
  \item Escreve uma equação do plano que passa em $A(1, 1, 1)$ e é paralelo ao plano $x + 2y - z = 4$.
  \item Determina o ponto de interseção da reta $(x, y, z) = (1, 0, 2) + t(1, 1, -1)$ com o plano $x + y + z = 6$.
  \item Verifica se os planos $2x - y + z = 1$ e $x + y - z = 3$ são perpendiculares.
  \item Escreve uma equação da reta perpendicular ao plano $3x - y + 2z = 5$ que passa na origem.
\end{exercicios}

\secao{Baricentro e Circuncentro}{Resolve.}
\begin{exercicios}
  \item Determina o baricentro do triângulo de vértices $A(1, 2)$, $B(5, 0)$ e $C(3, 7)$.
  \item Determina o baricentro do triângulo de vértices $A(1, 0, 2)$, $B(3, 3, 1)$ e $C(2, 0, 3)$.
  \item Determina o circuncentro do triângulo de vértices $A(0, 0)$, $B(6, 0)$ e $C(0, 4)$.
  \item Determina o circuncentro do triângulo de vértices $A(1, 1)$, $B(5, 1)$ e $C(3, 5)$.
  \item Mostra que, num triângulo retângulo, o circuncentro é o ponto médio da hipotenusa.
\end{exercicios}

\secao{Desafios}{Resolve os problemas.}
\begin{exercicios}
  \item Calcula a distância do ponto $P(1, 2, 3)$ ao plano $2x + y - 2z = 6$ (usa a projeção ortogonal de $P$ no plano).
  \item Determina a equação da reta de interseção dos planos $x + y + z = 3$ e $x - y = 1$.
  \item Mostra que os pontos $A(1, 2, 0)$, $B(3, 0, 2)$ e $C(4, -1, 3)$ são colineares.
  \item Num cubo de aresta 2 com um vértice na origem e arestas sobre os semieixos positivos, determina o ângulo entre duas diagonais espaciais.
  \item Determina a área do triângulo de vértices $A(0,0,0)$, $B(2,0,0)$ e $C(0,3,4)$.
\end{exercicios}

\end{document}
